
\subsection{Type 1: Independent Warping}

The most general approach to modeling warping is to take the \(n_w\) warping amplitudes \(\bm{\alpha}\) as independent configuration variables. 
This generalizes the approach of \cite{reissner1952nonuniform} for warping torsion.

\paragraph{Equilibrium}
Regardless of any kinematic assumptions on the primary kinematic variables \(\bm{x}\) and \(\FrameRotation\), the warping resultants \(\bm{w}\) and \(\bm{v}\) are governed by the same equilibrium equation:
\begin{equation}
\bm{w}' - \bm{v} + \bm{f}_{\mathrm{w}} = \mathbf{0}
\label{eq:bi-equilibrium}
\end{equation}
where $\bm{f}_{\mathrm{w}}$ is a generalized distributed load in the warping variables (typically taken to be zero). 
%
The complete system of equilibrium equations is
\begin{equation} 
\left\{\begin{array}{r}
(\FrameRotation \bm{n})^{\prime}+{\bm{f}}_n=\mathbf{0} \\
(\FrameRotation \bm{m})^{\prime}+\bm{x}^{\prime} \times (\FrameRotation \bm{n}) +{\bm{f}}_m=\mathbf{0} \\
\bm{w}^{\prime}-\bm{v} + \bm{f}_{\rm w} = \mathbf{0}
\end{array}\right. 
\label{eq:finite-equilibrium}
\end{equation}
where \(\bm{f}_n\) and \(\bm{f}_m\) are given body force and moment loadings, respectively. 
\citep{antman1974kirchhoffs,simo1991geometricallyexact,klinkel2003anisotropic}. 

% This is separated out of the equilibrium equation \eqref{eq:finite-equilibrium}, so that 
% all assumptions on and simplifications to the primary kinematic variables \(\bm{x}\) and \(\FrameRotation\) are contained within a \emph{primary} equilibrium operator \(\mathbb{B}_{\mathrm{p}}\). 


\vspace{1ex}
\begin{ambox}[Type 1: Independent Warping]{box:type1-warping}
A Type 1 constitutive response maps \(6+2n_w\) strains \((\bm{\gamma},\, \bm{\kappa},\, \bm{\alpha}',\, \bm{\alpha})\) to \(6+ 2n_w\) generalized stresses \((\bm{n},\, \bm{m},\,\bm{w},\bm{v})\)
\[
[\bm{n},\bm{m},\bm{w},\bm{v}] = \mathcal{S}(\bm{\gamma},\, \bm{\kappa},\, \bm{\alpha}',\, \bm{\alpha})
\]

\begin{equation}
\begin{aligned}
    \mathbb{B}_{\mathrm{p}} \bm{A}_*\begin{pmatrix} \bm{n} \\ \bm{m} \end{pmatrix} + \bm{f}_{\mathrm{p}} &= \dot{\bm{\mu}} \\ 
    \bm{w}' - \bm{v} + \bm{f}_{\mathrm{w}} &= \mathbf{0}
\end{aligned}
\end{equation}
\end{ambox}


\subsection{Type 1: Independent Warping}

The most general approach to model warping is to take the \(\nwarp\) warping amplitudes \(\bm{\alpha}\) as independent configuration variables. 
This follows the approach of \citep{reissner1952nonuniform} for warping torsion. 
Regardless of any kinematic assumptions on the primary kinematic variables \(\TraceLocation\) and \(\FrameRotation\), the warping resultants \(\bimoment\) and \(\bishear\) are governed by the same equilibrium equation:
\begin{equation*}
    \bimoment' - \bishear + \bm{f}_{\mathrm{w}} = \mathbf{0}
\end{equation*}
This is separated out of the equilibrium equation \eqref{eq:finite-equilibrium}, so that all assumptions on and simplifications to the primary kinematic variables \(\TraceLocation\) and \(\FrameRotation\) are contained within a \emph{primary} equilibrium operator \(\EquilibriumDerivative\). 

A \textbf{Type 1} constitutive formulation is completely general in the sense that it can readily return a response for all other formulation types.

\vspace{1ex}
\begin{ambox}[Type 1: Independent Warping]{box:type1-warping}
A Type 1 constitutive response maps \(6+2\, \nwarp\) strains \((\ShearStrain,\, \CurveStrain,\, \bm{\alpha}',\, \bm{\alpha})\) to \(6+ 2\, \nwarp\) generalized stresses \((\ShearResult,\, \CurveResult,\,\bm{w},\bm{v})\)
\[
[\bm{s}_{\mathrm{p}},\bm{w},\bm{v}] = \SectionEvaluation{\bm{e}_{\mathrm{p}},\, \bm{\alpha}',\, \bm{\alpha}}
\]
Equilibrium takes the form 
\begin{equation}
\begin{aligned}
    \mathbb{B}_{\rm p} \bm{\sigma}_{\mathrm{p}} + \bm{f}_{\mathrm{p}} &= \dot{\bm{\mu}} \\ 
    \bimoment' - \bishear + \bm{f}_{\mathrm{w}} &= \mathbf{0}
\end{aligned}
\end{equation}
\end{ambox}



\subsection{Type 2: Constrained Warping}

It is common to assume that the amplitudes of certain warping modes are proportional to a component of one of the primary strain variables \(\bm{e}_{\mathrm{p}} = (\bm{\gamma},\,\bm{\kappa})\). 
For example, torsional warping is often assumed to be proportional to the rate of twist \(\varkappa \triangleq \bm{\kappa}\cdot \FrameBasis\).
To formalize this, we partition the warping modes into three disjoint sets with indices \(\mathcal{I}_v\), \(\mathcal{I}_w\), and \(\mathcal{I}_f\) corresponding to algebraically constrained, differentially constrained, and free modes, respectively. 
We introduce projection operators \(\mathbf{L}_{\rm vv}, \mathbf{L}_{\rm ww} \in \mathbb{R}^{n_w \times n_w}\) onto the subspaces of algebraically and differentially constrained modes, satisfying the idempotency and orthogonality conditions:
\[
\mathbf{L}_{vv}^2 = \mathbf{L}_{vv}, \quad \mathbf{L}_{ww}^2 = \mathbf{L}_{ww}, \quad \mathbf{L}_{vv}\mathbf{L}_{ww} = \mathbf{L}_{ww}\mathbf{L}_{vv} = \mathbf{0}
\]

The constraints are expressed through linear maps \(\mathbf{L}_{\rm vp}, \mathbf{L}_{\rm wp} \in \mathbb{R}^{n_w \times 6}\) such that:
%
\begin{equation}
\left.
\begin{aligned}
\mathbf{L}_{\rm vv}\bm{\alpha} &= \mathbf{L}_{\rm vp}\bm{e}_{\mathrm{p}} \\
% \mathbf{L}_{\rm ww}\bm{\alpha}' &= \mathbf{L}_{\rm wp} \bm{e}_{\mathrm{p}}
\end{aligned}\right.
\label{eq:linear-warp-constraints}
\end{equation}
%
where the constraint maps satisfy the compatibility conditions \(\mathbf{L}_{\rm vv}\mathbf{L}_{\rm vp} = \mathbf{L}_{\rm vp}\) and \(\mathbf{L}_{\rm ww}\mathbf{L}_{\rm wp} = \mathbf{L}_{\rm wp}\). 
%
The warping amplitude vector is decomposed as:
\begin{equation}
\bm{\alpha} = \bar{\bm{\alpha}} + \mathbf{L}_{\rm vv}\bm{\alpha} + \mathbf{L}_{\rm ww}\bm{\alpha}
\label{eq:decompose-alpha}
\end{equation}
where the reduced warping vector \(\bar{\bm{\alpha}} \triangleq (\mathbf{I} - \mathbf{L}_{vv} - \mathbf{L}_{ww})\bm{\alpha}\) contains only the free components.

% \subsubsection{Algebraic Constraints}

\paragraph{Kinematics}
Using the constraints \eqref{eq:linear-warp-constraints}, the full strain vector can be expressed through the transformation:

\begin{equation}
\bm{e}=\begin{pmatrix}
\bm{\gamma} \\ \bm{\kappa} \\ \bm{\alpha}' \\\bm{\alpha}
\end{pmatrix}=\underbrace{\begin{bmatrix}
\mathbf{1} & \mathbf{0} & \mathbf{0} & \mathbf{0} \\ 
\mathbf{0} & \mathbf{1} & \mathbf{0} & \mathbf{0} \\ 
\mathbf{0} & \mathbf{0} & \mathbf{1} & \mathbf{0} \\
\mathbf{L}_{\rm vn} & \mathbf{L}_{\rm vm} & \mathbf{0} & \mathbf{1}-\mathbf{L}_{\rm vv}
\end{bmatrix}}_{\mathbf{L}}
\begin{pmatrix}
\bm{\gamma} \\ \bm{\kappa} \\ \bm{\alpha}' \\ \bar{\bm{\alpha}}
\end{pmatrix}
\label{eq:warping-projection}
\end{equation}
% \begin{equation}
% \bm{e}=\begin{pmatrix}
% \bm{\gamma} \\ \bm{\kappa} \\ \bm{\alpha}' \\\bm{\alpha}
% \end{pmatrix}=\underbrace{\begin{bmatrix}
% \mathbf{1} & \mathbf{0} & \mathbf{0} & \mathbf{0} \\ 
% \mathbf{0} & \mathbf{1}  & \mathbf{0} & \mathbf{0} \\ 
% \mathbf{L}_{\rm vn} &\mathbf{L}_{\rm vm} & \mathbf{1} - \mathbf{L}_{\rm ww} & \mathbf{0} \\
% \mathbf{L}_{\rm vn} & \mathbf{L}_{\rm vm} & \mathbf{0} & \mathbf{1}-\mathbf{L}_{\rm vv}
% \end{bmatrix}}_{\mathbf{L}}
% \begin{pmatrix}
% \bm{\gamma} \\ \bm{\kappa} \\ \bar{\bm{\alpha}}' \\ \bar{\bm{\alpha}}
% \end{pmatrix}
% \label{eq:warping-projection}
% \end{equation}
%
where \(\mathbf{L}_{\rm vn}\) and \(\mathbf{L}_{vm}\) are the partitions of \(\mathbf{L}_{\rm vp}\) corresponding to \(\bm{\gamma}\) and \(\bm{\kappa}\), respectively, and similarly for \(\mathbf{L}_{\rm wn}\) and \(\mathbf{L}_{\rm wm}\).
\begin{remark}
Here $\bar{\bm{\alpha}}$ partitions out the constrained amplitudes, and $\bm{\alpha}$ is the union of these constrained amplitudes with the unconstrained amplitudes. So, eg, given a two-mode warping model $\bm{\alpha} = (\alpha_1, \alpha_2)$, and one constraint, $\alpha_1 = \FrameBasis\cdot\bm{\kappa}$,
then \[
\mathbf{L}_{\rm vv} = \begin{bmatrix}
1 & 0 \\ 0 & 0
\end{bmatrix},
\qquad
\bar{\bm{\alpha}} = \bm{\alpha} - \mathbf{L}_{\rm vv}\bm{\alpha} = (0, \alpha_2),
\quad 
\bm{\alpha} = \bar{\bm{\alpha}} + (\alpha_1, 0)
\]
\end{remark}

\paragraph{Equilibrium}

To derive the equilibrium equations for the constrained system, we consider the Galerkin projection of the equilibrium residual. The weak form contribution from warping is:
\begin{equation*}
\mathcal{G}_{\chi}[\bm{\eta}] = \left< \mathbb{B}_{\mathrm{p}} \begin{pmatrix} \bm{n} \\ \bm{m} \end{pmatrix} ,\, \bm{\eta}_{\mathrm{p}} \right>
-\left< \bm{w}'-\bm{v},\, \delta\bm{\alpha} \right>
\end{equation*}

Using the decomposition \eqref{eq:decompose-alpha} and the constraint equations, the test functions \(\delta\bm{\alpha}\) can be expressed in terms of the reduced variables. 
After exploiting the linearity of the inner product and the transpose property of the constraint operators, we obtain:
\begin{equation*}
\mathcal{G}_{\chi}[\bm{\eta}] = 
\left<\begin{pmatrix} \bm{n} \\ \bm{m} \end{pmatrix} + \mathbf{L}_{\rm vp}^{\top}\bm{v} - \mathbf{L}_{\rm vp}^{\top}\bm{w}' ,\, \delta \bm{e}_{\mathrm{p}} \right>
-\left< \bm{w}'-\bm{v} ,\, \delta \bar{\bm{\alpha}} \right>
\end{equation*}

This motivates the definition of augmented stress resultants:
\begin{equation}
\begin{aligned}
\bar{\bm{n}} &\triangleq \bm{n} + \mathbf{L}_{\rm vn}^{\top}\bm{v} \\ % \\- \mathbf{L}_{\rm vn}^{\top}\bm{w}' \\
\bar{\bm{m}} &\triangleq \bm{m} + \mathbf{L}_{\rm vm}^{\top}\bm{v} \\ % \\- \mathbf{L}_{\rm vm}^{\top}\bm{w}' \\
\bar{\bm{w}} &\triangleq (\mathbf{1} - \mathbf{L}_{\rm vv} - \mathbf{L}_{\rm ww})^{\top}\bm{w} \\
\bar{\bm{v}} &\triangleq (\mathbf{1} - \mathbf{L}_{\rm vv} - \mathbf{L}_{\rm ww})^{\top}\bm{v}
\end{aligned}
\label{eq:augment-stress}
\end{equation}

and the corresponding resultant transformation \(\mathbf{L}^*\):
\begin{equation*}
\bar{\bm{s}}=\begin{pmatrix}
\bar{\bm{n}} \\ \bar{\bm{m}} \\ \bar{\bm{w}} \\ \bar{\bm{v}}
\end{pmatrix}=\underbrace{\begin{bmatrix}
\mathbf{1} & \mathbf{0} & \mathbf{0} & \mathbf{L}_{\rm vn}^{\top} \\ 
\mathbf{0} & \mathbf{1}  & \mathbf{0} & \mathbf{L}_{\rm vm}^{\top} \\ 
\mathbf{0} & \mathbf{0} & \mathbf{1} - \mathbf{L}_{\rm ww}^{\top} & \mathbf{0} \\
\mathbf{0} & \mathbf{0} & \mathbf{0} & \mathbf{1}-\mathbf{L}_{\rm vv}^{\top}
\end{bmatrix}}_{\mathbf{L}^*}
\begin{pmatrix}
\bm{n} \\ \bm{m} \\ \bm{w} \\ \bm{v}
\end{pmatrix}
\end{equation*}
%
\vspace{1ex}
\begin{ambox}[Type 2: Constrained Warping]{box:warp2}
For linear constraints on warping with the form of \eqref{eq:linear-warp-constraints}, the Type 1 constitutive response is transformed through:
\[
\bar{\bm{s}} = \mathbf{L}^{\top}\mathcal{S}(\mathbf{L}\bar{\bm{e}})
\]
where \(\bar{\bm{s}} = (\bar{\bm{n}}, \bar{\bm{m}}, \bar{\bm{w}}, \bar{\bm{v}})\) and \(\bar{\bm{e}} = (\bm{\gamma}, \bm{\kappa}, \bar{\bm{\alpha}}', \bar{\bm{\alpha}})\).
The momentum balance becomes:
% \begin{equation}
% \begin{aligned}
%     \mathbb{B}_{\mathrm{p}} \begin{pmatrix} \bar{\bm{n}} \\ \bar{\bm{m}} \end{pmatrix} + \bm{f}_{\mathrm{p}} &= \dot{\bm{\mu}} \\ 
%     \bar{\bm{w}}' - \bar{\bm{v}} + \bar{\bm{f}}_{\mathrm{w}} &= \mathbf{0}
% \end{aligned}
% \label{eq:type2-equilibrium}
% \end{equation}

\begin{equation}
\begin{aligned}
    \mathbb{B}_{\mathrm{p}} \bm{A}_* \left[\bar{\bm{s}}_{\mathrm{p}} + \mathbf{L}_{\rm vp}^{\top}\bm{w}'\right]+ \bm{f}_{\mathrm{p}} &= \dot{\bm{\mu}} \\ 
    \bar{\bimoment}' - \bar{\bishear} + \bm{f}_{\mathrm{w}} &= \mathbf{0}
\end{aligned}
\label{eq:type2-equilibrium}
\end{equation}
\end{ambox}

\begin{remark}
The treatment of the $\mathbf{L}_{\rm vp}^{\top}\bm{w}'$ term in \eqref{eq:type2-equilibrium} reveals a fundamental distinction between section-level and element-level operations. The constitutive evaluation at a section operates pointwise and cannot access spatial derivatives. Therefore:
\begin{itemize}
\item The \textbf{section wrapper} receives $(\bm{\gamma}, \bm{\kappa}, \bar{\bm{\alpha}}', \bar{\bm{\alpha}})$, expands them via $\mathbf{L}$ to obtain the full strain $\bm{e}$, evaluates the Type 1 constitutive response $\mathcal{S}(\bm{e})$, and returns the partially augmented resultants:
\[
\bar{\bm{n}} = \bm{n} + \mathbf{L}_{\rm vn}^{\top}\bm{v}, \quad 
\bar{\bm{m}} = \bm{m} + \mathbf{L}_{\rm vm}^{\top}\bm{v}, \quad
\text{and the full } \bm{w}
\]
\item The \emph{element} must then compute $\bm{w}'$ through differentiation of its interpolated bimoment field and apply the correction $-\mathbf{L}_{\rm vp}^{\top}\bm{w}'$ to the primary equilibrium.
\end{itemize}
This division implies that Type 2 elements require higher-order interpolation to accurately compute $\bm{w}'$, making them more complex to implement than standard beam elements. 
This motivates the Type 3 formulation, where the assumption $\bar{\bm{\alpha}}' = \mathbf{0}$ (uniform warping) eliminates the $\bm{w}'$ term entirely, allowing all operations to remain at the section level and enabling the use of standard 6-DOF beam elements without modification.
\end{remark}

\subsection{Type 3: Uniform Warping}

For algebraic constraints in Type 3, we seek to derive a case that is appropriate for elements that cannot supply $\bm{\alpha}'$ to the section. 


Type 3 imposes two conditions on the constrained warping modes:
\begin{equation}
\mathbf{L}_{\rm vp}\bm{e}_{\rm p}' = \mathbf{0} \quad \text{and} \quad \langle\bm{w}' - \bm{v}, \mathbf{L}_{\rm vp}\delta\bm{e}_{\rm p}\rangle = \mathbf{0}
\label{eq:type3-conditions}
\end{equation}
which is consistent with the assumption that warping occurs with a constant amplitude along the length of the body. 
The first enforces uniform constrained warping amplitudes, while the second ensures that bimoment gradients produce no work through the constrained modes. Together, these eliminate the $\mathbf{L}_{\rm vp}^T\bm{w}'$ term from the equilibrium.

\vspace{1ex}
\begin{ambox}[Type 3: Uniform Warping]{box:warp3}
Under conditions \eqref{eq:type3-conditions}, the constitutive response and equilibrium simplify to:
\begin{equation*}
[\bm{n}, \bm{m}] = \mathcal{S}_{\rm wrapped}(\bm{\gamma}, \bm{\kappa}, \bar{\bm{\alpha}}', \mathbf{L}_{\rm vp}\bm{e}_{\rm p} + \bar{\bm{\alpha}})
\end{equation*}
with momentum balance:
\begin{equation}
% \mathbb{B}_{\rm p} \begin{pmatrix} \bm{n} \\ \bm{m} \end{pmatrix} + \bm{f}_{\rm p} = \dot{\bm{\mu}}
\mathbb{B}_{\rm p} \bm{A}_* \begin{pmatrix} \bm{n} \\ \bm{m} \end{pmatrix} + \bm{f}_{\mathrm{p}} = \dot{\bm{\mu}}
\end{equation}
Standard 6-DOF beam elements can thus be used without modification, as they need only receive the primary stress resultants $(\bm{n}, \bm{m})$.
\end{ambox}

 This model is not able to enforce equilibrium of the constrained bishears/bimoments, but rather relies on the assumption that these naturally equilibrate. 


\subsection{Type 4: Neglected Warping}
\[
\mathbf{L}_{\mathrm{vp}} \bm{\alpha} =\mathbf{L}_{\mathrm{vp}} \bm{\alpha}^{\prime}=\mathbf{0}
\]
\begin{ambox}[Type 4: Neglected Warping]{box:warp4}
Balance of momentum is stated purely in terms of \(\bm{s}_{\mathrm{p}}\)
\begin{equation}
\begin{aligned}
    \mathbb{B}_{\rm p} \bm{A}_* \bm{s}_{\mathrm{p}}+ \bm{f}_{\mathrm{p}} &= \dot{\bm{\mu}}
\end{aligned}
\end{equation}
\end{ambox}


\subsection{Summary}

A library of finite elements typically contains ``Type 1/2", and ``Type 3" models. 
Type 1/2 elements have over 6 dofs per node and are somewhat less standard. 
Type 3 elements are classical with 6 dofs/node and expect only $(\bm{n},\bm{m})$ from the section. 
We want one section implementation which supports all element types. 
The section always recieves \((\gamma,\kappa,\alpha',\alpha)\) and evaluates the full $\bm{s} = (\bm{n},\bm{m},\bm{w},\bishear)$, and a ``wrapper" interface uses the $\mathbf{L}$ operators that handles the transformation $\bar{\bm e}\rightarrow\bm{e}$ and $\bm{s}\rightarrow\bar{\bm s}$. 
A Type 2 element notably may require higher order interpolations because derivatives of \(\bm{\gamma}\) and \(\bm{\kappa}\) will be required for \(\mathbf{L}_{\rm vv}\bm{\alpha}'\) components that are constrained (eg, the Timoshenko/Vlasov non-uniform torsion formulation).
A Type 3 element may be regarded as one where these higher order terms are neglected, and the derivatives of these constrained amplitudes are set to zero \(\mathbf{L}_{\rm vv}\bm{\alpha}' = \mathbf{0}\).
Thus, these elements only implement the $\mathbb{B}_{\rm p}$ partition of the full equilibrium operator. 
We assume that type 3 elements take 6-DOF nodal values, interpolate to find $\bm\gamma,\bm\kappa$, at integration points along $\reflenx$, then pass these to the section ``wrapper" which forms the full $\bm{e} = (\bm\gamma,\bm\kappa,\bm\alpha,\bm\alpha')$, passes it to the section in exchange for $\bm{s} = (\bm{n},\bm{m},\bm{w},\bishear)$, then condenses this to return 6 appropriate resultants to the element. 
The section wrapper and section itself can know nothing about where they lie along $\reflenx$. 
They only know the $\mathbf{L}$ operators and pointwise values of $\bm{e},\bm{s}$.

